Utilizando arquiteturas de backup, os sistemas críticos possuem redundância em hardware e software para lidar com falhas de componentes.

\subsection{Redundância nos sistemas espaciais }

Os estudos de \cite{sklaroff1976} dizem que o uso de sistemas redundantes em ambientes críticos ou de microgravidade tem sido utilizados para garantir confiabilidade e tolerância a falhas. Quando se trata de sistemas de controle de voo, como os utilizados no ônibus espacial (\textit{Space Shuttle}), por exemplo, a redundância é utilizada para garantir que funções essenciais sejam mantidas mesmo diante de falhas. O conjunto de computadores do ônibus espacial era composto por cinco unidades idênticas, sendo que quatro deles operavam de forma redundante para funções críticas de voo, enquanto que a quinta realizava tarefas não importantes. Durante as missões em que a recuperação de falhas deveria ocorrer em frações de segundo, a redundância permitia que comparações entre as saídas dos computadores identificassem possíveis falhas e isolassem unidades defeituosas, garantindo a continuidade da operação.

A arquitetura dos computadores utilizavam um esquema de votação para validar os comandos e detectar as falhas. O funcionamento se dava da seguinte forma: quando dois computadores apresentavam resultados divergentes em relação aos demais, esses eram considerados defeituosos, e suas saídas ou resultados eram ignorados. O gerenciador de redundâncias incluía a capacidade de detectar e identificar falhas antes do lançamento ou durante a missão, no caso de falhas múltiplas, os computadores restantes utilizavam o autodiagnóstico para avaliar a continuidade operacional do sistema. Esse modelo de redundância proporcionava uma cobertura majoritária de falhas, reduzindo a probabilidade de perda do sistema devido a mau funcionamento de componentes isolados.

A arquitetura computacional do ônibus espacial foi baseada em um conjunto de computadores IBM AP-101 \citep{ibm_space_shuttle}. Cada um continha uma unidade central de processamento e um processador de entrada e saída, conectados a barramentos seriais de comunicação. Desses, quatro eram computadores que executavam as mesmas funções de controle de voo, enquanto que, o último quinto era programado para funções auxiliares. Os subsistemas do veículo espacial eram distribuídos entre barramentos redundantes, de forma que cada barramento possuía canais dedicados para comunicação com diferentes computadores. Isso permitia que cada unidade recebesse os mesmos dados de entrada e produzisse comandos idênticos para os atuadores e sensores de bordo.

A redundância nos sistemas de controle não se limitava ao hardware computacional, mas também envolvia todos os sensores, atuadores e barramentos de dados. Alguns subsistemas críticos, como atuadores das superfícies de controle e motores principais, operavam em canais redundantes, do qual comandos independentes eram gerados e comparados antes de serem executados. Tal redundância evitava que falhas isoladas comprometessem a capacidade de controle da espaçonave. Além disso, o uso de algoritmos de detecção de falhas permitia que comandos incoerentes fossem identificados e desconsiderados antes de afetar a operação da missão.

Além disso, a gestão da redundância também incluía a sincronização de processos entre os computadores, como diferentes unidades podiam operar com baixas variações de tempo na execução de cálculos críticos, e era necessário garantir que todos os computadores sincronizassem seus estados periodicamente para evitar desvios no comportamento do sistema. Esse processo de sincronização era realizado por meio de sinais discretos e por um software especializado que coordenava a execução dos programas entre as unidades redundantes. Portanto, evitava-se que pequenas diferenças nos tempos de processamento levassem a falhas na coordenação entre os computadores do sistema.

A definição de requisitos para o gerenciamento da redundância considerava critérios como simplicidade de hardware, flexibilidade operacional e baixa utilização de recursos computacionais. O principal requisito era a capacidade de identificar falhas antes do lançamento e durante a missão, garantindo que computadores defeituosos fossem isolados de forma eficiente. Era essencial que falhas em um computador não resultassem na geração de comandos incorretos ou na identificação errônea de outras unidades como defeituosas. Para atender a esses requisitos, foram desenvolvidos mecanismos que permitiam aos computadores do conjunto redundante avaliar continuamente o desempenho das demais unidades e ajustar o funcionamento do sistema conforme necessário.

A redundância aplicada ilustra um modelo de gestão de falhas para sistemas críticos, combinando hardware redundante, algoritmos de detecção de falhas e técnicas de sincronização de processos. Esse modelo proporcionava um nível de confiabilidade elevado, pois reduzia os riscos de falhas catastróficas durante as missões espaciais. A experiência adquirida com essa abordagem contribuiu para o desenvolvimento de novas arquiteturas de sistemas tolerantes a falhas, influenciando projetos posteriores de veículos espaciais e sistemas aeronáuticos \citep{sklaroff1976}.
